\documentclass[11pt]{article}
\usepackage[margin=1in]{geometry}
\usepackage{hyperref}
\usepackage{listings}
\usepackage{xcolor}
\usepackage{enumitem}

\lstset{
    basicstyle=\ttfamily\small,
    backgroundcolor=\color{gray!10},
    frame=single,
    framerule=0.5pt,
    breaklines=true,
    showstringspaces=false,
    language=bash,
    morekeywords={sudo,apt,cmake,ninja,git,cd,mkdir,export}
}

\title{Building femm-qt on Linux}
\author{D.\ Meeker}
\date{\today}

\begin{document}
\maketitle

\section{Overview}

femm-qt is a cross-platform port of FEMM (Finite Element Method Magnetics)
using Qt~6 for the GUI and CMake/Ninja for the build system.  This document
describes how to build it from source on Linux.  The instructions target
Ubuntu 24.04~LTS but should work with minor changes on other distributions
(Debian, Fedora, Arch, etc.).

The build produces six executables:
\begin{itemize}[nosep]
  \item \texttt{femm\_gui} --- the main GUI application
  \item \texttt{femm\_headless} --- headless IPC server (no GUI)
  \item \texttt{fkn} --- magnetics solver
  \item \texttt{belasolv} --- electrostatics solver
  \item \texttt{hsolv} --- heat flow solver
  \item \texttt{csolv} --- current flow solver
\end{itemize}

\section{Prerequisites}

\subsection{Build Tools}

Install the compiler toolchain, CMake, Ninja, and git:
\begin{lstlisting}
sudo apt update
sudo apt install build-essential cmake ninja-build git
\end{lstlisting}
This provides \texttt{g++}, \texttt{cmake}, \texttt{ninja}, and \texttt{git}.

\subsection{Qt 6}

Install the Qt~6 development libraries:
\begin{lstlisting}
sudo apt install qt6-base-dev libqt6widgets6
\end{lstlisting}
On Ubuntu 24.04, this installs Qt~6.4 or later from the standard
repositories.  No separate Qt installer or online account is needed.

To verify the installation:
\begin{lstlisting}
qmake6 --version
\end{lstlisting}

\subsection{Other Distributions}

On Fedora:
\begin{lstlisting}
sudo dnf install gcc-c++ cmake ninja-build git qt6-qtbase-devel
\end{lstlisting}

On Arch Linux:
\begin{lstlisting}
sudo pacman -S base-devel cmake ninja git qt6-base
\end{lstlisting}

\section{Getting the Source}

If you have a source zip, unzip it:
\begin{lstlisting}
unzip femm-qt-src.zip
cd femm-qt
\end{lstlisting}

Or, if the source is hosted on GitHub, clone the repository:
\begin{lstlisting}
git clone <repository-url> femm-qt
cd femm-qt
\end{lstlisting}

\section{Building}

From the \texttt{femm-qt} directory, run the build script:
\begin{lstlisting}
./linux_build.sh
\end{lstlisting}

This checks that the build tools are present, configures the project
with CMake, and builds it with Ninja.  All executables are placed in
\texttt{build/bin/}.

To wipe the build directory and rebuild from scratch, pass
\texttt{clean}:
\begin{lstlisting}
./linux_build.sh clean
\end{lstlisting}

If the build fails because CMake cannot find Qt~6, see
\textit{Manual configuration} under Troubleshooting.

\section{Installing}

You can run femm-qt from \texttt{build/bin} indefinitely; nothing requires
an install.  To put it somewhere permanent instead:
\begin{lstlisting}
./linux_build.sh install
\end{lstlisting}

This installs what is already built --- it does not rebuild --- into your
home directory:
\begin{lstlisting}
~/.local/opt/femm-qt/   the program, its solvers, the material
                        libraries, the examples and the bindings
~/.local/bin/           femm_gui, femm_headless, femm-lua
~/.local/share/         menu entry and icon
\end{lstlisting}

\texttt{\textasciitilde/.local/bin} is already on \texttt{PATH} on current
distributions, so \texttt{femm\_gui} then works from anywhere and FEMM
appears in the application menu.  Pass a different prefix if you would
rather it went elsewhere:
\begin{lstlisting}
./linux_build.sh install /opt/femm
\end{lstlisting}

It installs into your home rather than \texttt{/usr/local} on purpose.
FEMM finds its solvers, scripting assets and material libraries beside
its own executable, so a single self-contained directory needs no
special handling --- and, more importantly, the material libraries and
the examples stay \emph{writable}.  Both need to be: the material editor
saves the libraries, and every example script opens a model and then
immediately saves it under another name.  In a system directory they
would be read-only and neither would work.

To uninstall, delete the three paths listed above; the installer prints
the exact commands when it finishes.

\section{Running}

Launch the GUI:
\begin{lstlisting}
./bin/femm_gui
\end{lstlisting}

Or the headless server:
\begin{lstlisting}
./bin/femm_headless
\end{lstlisting}


If you installed as above, use the copies on your \texttt{PATH} instead,
from any directory:
\begin{lstlisting}
femm_gui
femm_headless
\end{lstlisting}

\subsection{Setting FEMM\_DIR}

The language bindings (OctaveFEMM, MathFEMM, pyFEMM, JuliaFEMM) locate
the FEMM executables via the \texttt{FEMM\_DIR} environment variable.
Set it to the \texttt{bin} directory:
\begin{lstlisting}
export FEMM_DIR=/path/to/femm-qt/build/bin
\end{lstlisting}
or, if you installed it, at the program directory:
\begin{lstlisting}
export FEMM_DIR=~/.local/opt/femm-qt
\end{lstlisting}
Add this to your \texttt{\textasciitilde/.bashrc} to make it persistent.
The binding files themselves are in \texttt{bindings/} under whichever of
those you use --- \texttt{mfiles} for OctaveFEMM, and \texttt{mathfemm},
\texttt{scifemm}, \texttt{juliafemm}, \texttt{pyfemm} and
\texttt{luafemm} beside it.

\subsection{Desktop Integration (Optional)}

\emph{The \texttt{install} step above does this for you, with an icon.}
The recipe below is for running in place out of \texttt{build/bin}.

To add femm-qt to your application menu, create a \texttt{.desktop} file:
\begin{lstlisting}
cat > ~/.local/share/applications/femm-qt.desktop << 'EOF'
[Desktop Entry]
Name=FEMM
Comment=Finite Element Method Magnetics
Exec=/path/to/femm-qt/build/bin/femm_gui
Type=Application
Categories=Science;Engineering;
EOF
\end{lstlisting}

\section{Using with MATLAB or Octave}

\subsection{Octave}

Install Octave and its Java support:
\begin{lstlisting}
sudo apt install octave default-jdk
\end{lstlisting}
Note that in Ubuntu, the version of Octave installed by App Center does \textit{not} have Java support, so the installation by \texttt{apt}, which does have Java support, is needed.
Find the path to Java by typing:
\begin{lstlisting}
readlink -f $(which java)
\end{lstlisting}
at the Bash command line.  This usually returns something like: 
\begin{lstlisting}
/usr/lib/jvm/java-21-openjdk-amd64/bin/java
\end{lstlisting}
Edit your \verb+~/.bashrc+ file and add as the last lines:
\begin{lstlisting}
export JAVA_HOME=/usr/lib/jvm/java-21-openjdk-amd64
export FEMM_DIR=/path/to/femm-qt/build/bin
\end{lstlisting}
Run Octave and confirm that Java is running by typing
\begin{lstlisting}
version -java
\end{lstlisting}
at the Octave prompt.  It should return something like: 
\\ \\
\begin{small} \texttt{ans =  Java 21.0.10+7-Ubuntu-124.04 with Ubuntu OpenJDK 64-Bit Server VM mixed mode, sharing} \end{small} \\
which indicates that Java is properly installed.  Then add the path to the the OctaveFEMM mfiles to Octave by typing the following at the Octave command line:
\begin{lstlisting}
addpath("/path/to/femm-qt/octavefemm")
savepath
\end{lstlisting}
Now, the \texttt{openfemm} function in the \texttt{octavefemm} directory should find and launch \texttt{femm\_gui}
or \texttt{femm\_headless} automatically.

\subsection{MATLAB}

Set \texttt{FEMM\_DIR} before launching MATLAB, or set it within MATLAB:
\begin{lstlisting}
setenv('FEMM_DIR', '/path/to/femm-qt/build/bin')
\end{lstlisting}
MATLAB on Linux uses its bundled JRE for socket communication with
femm-qt, which should work without additional configuration.  The path to the OctaveFEMM mfiles must also be set using the \verb+pathtool+ command.

\subsection{Python}

The Python interface is/will be available as pyfemm-qt.  Install it via:
\begin{lstlisting}
pip install pyfemm-qt
\end{lstlisting}
To install the version that comes with the source code distribution instead, change to the \verb+pyfemm+ directory in the source distribution and type:
\begin{lstlisting}
pip install -e .
\end{lstlisting}
On Ubuntu 24.04 and later, pip may refuse to install into the system
Python with an ``externally-managed-environment'' error.  Since pyfemm
is a single file with no dependencies, it is safe to override this:
\begin{lstlisting}
pip install -e . --break-system-packages
\end{lstlisting}

\section{Troubleshooting}

\subsection{Manual configuration}

The build script just wraps the standard CMake and Ninja steps.  To run
them by hand:
\begin{lstlisting}
mkdir build && cd build
cmake .. -G Ninja -DCMAKE_BUILD_TYPE=Release
ninja
\end{lstlisting}
CMake finds Qt~6 from the system packages automatically.  If it does
not, pass the Qt prefix explicitly:
\begin{lstlisting}
cmake .. -G Ninja \
    -DCMAKE_BUILD_TYPE=Release \
    -DCMAKE_PREFIX_PATH=/usr/lib/x86_64-linux-gnu/cmake/Qt6
\end{lstlisting}

\subsection{Qt not found}

If CMake reports that it cannot find Qt6, check that the development
package is installed:
\begin{lstlisting}
dpkg -l | grep qt6-base-dev
\end{lstlisting}
If installed but not found, pass \texttt{-DCMAKE\_PREFIX\_PATH} pointing
to the Qt~6 CMake config directory, typically
\texttt{/usr/lib/x86\_64-linux-gnu/cmake/Qt6}.

\subsection{Missing OpenGL headers}

On minimal server installations or VMs without a desktop environment,
Qt may require OpenGL development headers:
\begin{lstlisting}
sudo apt install libgl1-mesa-dev libglu1-mesa-dev
\end{lstlisting}

\subsection{VirtualBox Guest Additions}

If running in a VirtualBox VM, install Guest Additions for proper display
scaling and shared folders:
\begin{lstlisting}
sudo apt install virtualbox-guest-utils \
    virtualbox-guest-x11
\end{lstlisting}
Reboot after installation.  For shared folders, add your user to the
\texttt{vboxsf} group:
\begin{lstlisting}
sudo usermod -aG vboxsf $USER
\end{lstlisting}

\subsection{MESA / OpenGL errors in VirtualBox}

VirtualBox's virtual GPU may not support the OpenGL version Qt expects,
producing errors like ``ZINK: failed to choose pdev'' or ``failed to
create dri2 screen.''  Two fixes:

\begin{enumerate}[nosep]
\item Enable 3D Acceleration in VirtualBox (Settings $\rightarrow$
      Display $\rightarrow$ check ``Enable 3D Acceleration'' and move the slider to at least 128MB video memory out in Expert mode), or
\item Force Qt to use software rendering by adding
\begin{lstlisting}
export LIBGL_ALWAYS_SOFTWARE=1
\end{lstlisting}
to \verb+~/.bashrc+
\end{enumerate}

\end{document}
